\chapter{What Probability Really Is}

As promised at the close of the preface, we start with a coin in your hand. You
flip it. While it is in the air, you cannot predict whether this flip will land
heads or tails, yet you would gladly bet on heads at even odds and refuse ten-to-one
against. Somehow, in the middle of total ignorance about \emph{this} flip, you have a
confident, quantitative opinion. You can put a \emph{number} on your uncertainty.

That number is a probability. This whole book is about where such numbers come
from, what they mean, why they obey the rules they do, and how far they reach --
from card tables in Renaissance Italy to spam filters in your inbox, from insurance
premiums to quantum mechanics and clinical trials. But before we compute a single
thing, we have to answer a deceptively hard question: what \emph{is} a probability,
really?

This chapter is the foundation, so we will go slowly and take the idea apart to
see what makes it tick.

\section{Two Good Answers to ``What Does 1/2 Mean?''}
\label{ch01:two-answers}

Ask ten people what ``the probability of heads is $1/2$'' means and you will get,
mostly, two genuinely different answers. Both are correct. They are correct about
\emph{different things}, and confusing them is the source of an enormous amount of
muddle. Both deserve a hearing.

\subsection{The frequentist answer: probability is long-run frequency}
\label{ch01:frequentist}

The first answer points outward, at the world. ``Probability $1/2$'' means: if you
flipped this coin a great many times, about half of the flips would be heads. The
probability is the \emph{limiting relative frequency} of the outcome in a long
run of repetitions.

This is the interpretation most people absorb in school, and it has a bracing,
no-nonsense appeal. It says probability is not in your head; it is a measurable
property of a repeatable experiment, like the boiling point of water. You do not
have to argue about it -- you go flip the coin ten thousand times and count.

\begin{intuitionbox}
Imagine a long ledger. Each row is one flip: H or T. Walk down the ledger and, at
every row, write down the fraction of H's \emph{so far}. Early on that fraction
lurches around -- one flip in, it is either $0$ or $1$; ten flips in it might be
$0.3$ or $0.7$. But keep walking. The fraction settles. It homes in on a number
and stops wandering far from it. The frequentist says: \emph{that number} is the
probability. It is the address the running fraction is heading toward.
\end{intuitionbox}

We can watch this happen. Figure~\ref{ch01:fig-stabilize} shows the running
fraction of heads for two simulated coins, one fair ($p = 0.5$) and one biased
($p = 0.7$), as the number of flips grows. Early on, the running fraction swings widely;
as the number of flips grows, it settles toward the true value.

\begin{figure}[H]
\centering
\begin{tikzpicture}
\begin{axis}[
  width=0.92\textwidth, height=6.2cm,
  title style={font=\small},
  title={Running frequency of heads: stabilizing toward true probability},
  xlabel={number of flips $n$}, ylabel={fraction of heads so far},
  xmin=1, xmax=500, ymin=0, ymax=1,
  xmode=log, log basis x=10,
  grid=both, grid style={gray!18},
  legend pos=north east, legend cell align=left,
]
% fair coin running fraction (precomputed, smooth-ish toward 0.5)
\addplot[sourceblue, thick] coordinates {
(1,1)(2,0.5)(3,0.667)(5,0.6)(8,0.5)(12,0.583)(20,0.55)(30,0.467)
(50,0.52)(80,0.488)(120,0.508)(200,0.49)(300,0.503)(400,0.4975)(500,0.502)
};
\addlegendentry{fair coin, $p=0.5$}
% biased coin toward 0.7
\addplot[densityred, thick] coordinates {
(1,1)(2,1)(3,0.667)(5,0.8)(8,0.75)(12,0.667)(20,0.75)(30,0.7)
(50,0.72)(80,0.6875)(120,0.708)(200,0.695)(300,0.703)(400,0.6975)(500,0.701)
};
\addlegendentry{biased coin, $p=0.7$}
\draw[gray, dashed] (axis cs:1,0.5) -- (axis cs:500,0.5);
\draw[gray, dashed] (axis cs:1,0.7) -- (axis cs:500,0.7);
\end{axis}
\end{tikzpicture}
\caption{Running fraction of heads versus number of flips (horizontal axis on a
log scale, so the early flips are not crushed against the left edge). Both curves
fluctuate at first and then settle toward their true probabilities (dashed
lines); the frequentist defines probability as this limiting value.}
\label{ch01:fig-stabilize}
\end{figure}

\begin{trapbox}
The frequentist definition has a well-known limitation worth addressing
immediately. The phrase ``in the long run'' hides an
infinity. You can never actually flip a coin infinitely many times, so the
limiting frequency is never literally observed -- it is an idealization. Worse, it
is circular if you are not careful: to say the fraction ``converges'' is itself a
statement that only holds \emph{with probability one}, so we have used the word
probability to define probability. We will untangle this when we lay down
Kolmogorov's axioms in Chapter~3, and again when we prove the Law of Large
Numbers later in the book. For now, hold the frequentist picture as an
\emph{interpretation}, a way of attaching meaning to the math -- not as the math
itself.
\end{trapbox}

The frequentist view is at its most natural when the experiment really is
repeatable under stable conditions: dice, cards, manufactured parts coming off a
line, radioactive atoms, the click of a Geiger counter. In all these cases ``do it
again'' means something concrete.

\subsection{The Bayesian answer: probability is degree of belief}
\label{ch01:bayesian}

Now consider a different kind of question. \emph{What is the probability that it
rained in Rome on the day Julius Caesar was assassinated?} Or: \emph{what is the
probability that this particular patient has the disease, given this particular
test result?} Or: \emph{what is the probability that the missing airliner is in
this search box?}

There is no ``long run'' here. Caesar was assassinated exactly once. This patient
is one person. We cannot re-run history under stable conditions and count. The
frequentist definition, taken literally, simply has nothing to say -- the event
either happened or it did not; the frequency is $0$ or $1$ and we just do not know
which.

And yet these are exactly the situations where we most desperately need
probability. So there is a second answer, older in spirit and in some ways deeper.
\emph{Probability is a measure of how strongly a rational agent believes
something, given the information they have.} It lives in the mind, not in the
coin. ``Probability $0.7$'' means: I am willing to bet on this at odds of
$7$-to-$3$; I would be indifferent between winning a prize if the event happens
and winning that same prize if I draw a green ball from an urn that is $70\%$
green.

\begin{intuitionbox}
A degree of belief is not a feeling, not a vague hunch. It is pinned down by
\emph{action}. Suppose I offer you a ticket that pays \$1 if it rained in Rome on
the Ides of March, 44 BC, and \$0 otherwise. How much would you pay for that
ticket? If the most you would pay is 30 cents, then -- on pain of being a sure
loser to a clever bookmaker -- your degree of belief in ``it rained'' is $0.3$.
Probability, on this view, is the price of a bet, made coherent.
\end{intuitionbox}

The Bayesian (named for Thomas Bayes, an 18th-century minister whose posthumous
essay we will study in depth) takes degree of belief as primary and then insists
on \emph{coherence}: your various beliefs must fit together without contradiction.
An important fact -- which we will prove later in the book via so-called
``Dutch book'' arguments -- is that the only way to keep your betting prices from
being exploitable is to make them obey the exact same rules of probability that the
frequentist derives from counting. Two completely different starting points --
pricing bets coherently and counting outcomes in the long run -- lead to the
exact same algebra.

\begin{marginnotebox}
The word ``probable'' long predates the mathematics. In medieval and early modern
usage, an opinion was \emph{probabilis} if it was approvable, worthy of belief --
backed by authority or argument. (Echoes survive in ``probity'' and ``approve.'')
The shift from probability-as-credibility to probability-as-number was gradual,
and the two senses, belief and frequency, were tangled together for centuries
before anyone realized they could come apart. When you feel the pull of both
interpretations at once, you are feeling a genuine three-hundred-year-old
ambiguity, not a personal confusion.
\end{marginnotebox}

\subsection{When each is natural -- and why the fight is overblown}
\label{ch01:which-when}

Generations of statisticians waged something close to a holy war over which
interpretation is ``correct.'' You can largely sidestep it. Here is the practical
truth.

\begin{itemize}
\item Use the \textbf{frequentist} picture when the experiment is genuinely
repeatable and you care about long-run behavior: quality control, casino odds,
estimating an error rate, the performance of a procedure over many future cases.
\item Use the \textbf{Bayesian} picture when you face a \emph{one-off} uncertainty
and want to combine evidence with prior knowledge: medical diagnosis for this
patient, the location of one lost ship, whether one hypothesis is true given one
dataset.
\end{itemize}

\begin{aestheticbox}
The key point to keep in mind is that \emph{the mathematics is the same either way}. Both camps agree that
probabilities are numbers between $0$ and $1$, that the probability of ``A or B''
for incompatible events is the sum, that conditional probability multiplies along
a tree. The interpretations differ only in what those numbers \emph{mean} and
hence in what questions they let you ask. So we can build the entire machine --
the axioms, the random variables, the distributions, the limit theorems -- on
neutral ground, and then plug in whichever interpretation fits the problem in
front of you. That is exactly what we will do. The philosophy is real and we will
honor it, but it never blocks the calculation.
\end{aestheticbox}

\section{Why Uncertainty Needs Mathematics}
\label{ch01:why-math}

You might reasonably ask: if probability is about not knowing, why do we need
\emph{mathematics}? Mathematics is the science of certainty -- of proof, of $2+2=4$
forever. What business does it have with the unknown?

The answer is that raw intuition makes systematic mistakes whenever chance is
involved. Our gut serves us well in everyday life, but it misfires badly on
probabilities. We see faces in clouds and patterns in lottery numbers.
We think a basketball player who hit five shots is ``hot'' when the streak is well
within randomness. We are terrified of shark attacks and blase about car rides.
We will spend a whole chapter on these cognitive traps, but here is the headline:
\emph{quantitative uncertainty does not yield to common sense; it yields to
calculation.}

\begin{intuitionbox}
Consider a quick gut-check (the third puzzle from the preface). A test for a
rare disease is ``$99\%$ accurate.'' You test positive. What is the probability
you have the disease? Most people, including many doctors, blurt out ``about
$99\%$.'' The real answer, when the disease is rare, can be under $10\%$. Your
intuition is not slightly off here; it is off by a factor of ten. Only
arithmetic -- which we will do carefully when we reach Bayes' theorem in
Chapter~5 -- can rescue you. This single example has changed real medical
practice.
\end{intuitionbox}

Mathematics does three things for uncertainty that intuition cannot.

\textbf{First, it forces consistency.} If you believe $A$ has probability $0.6$,
then you are \emph{committed} to ``not $A$'' having probability $0.4$ -- you cannot
choose otherwise without contradicting yourself. The rules act like the rules of
accounting: they do not tell you the numbers, but once you set some, they pin down
others, and they catch you when your story does not add up.

\textbf{Second, it lets us combine uncertainties.} Real problems chain together --
the probability of a system failing depends on its parts, the probability of a
diagnosis depends on several tests. Intuition cannot reliably multiply and add a
dozen partial beliefs; a calculus can.

\textbf{Third, it reveals stable laws hiding inside the chaos.} Individual random
events are unpredictable, but \emph{aggregates of them are remarkably
predictable}. You cannot say whether the next customer buys, but a store can
forecast next month's sales to within a percent. You cannot predict one molecule,
but you can write the gas laws. Out of a swarm of unknowables, sharp regularities
crystallize. Probability is the mathematics of that crystallization.

\begin{happeningbox}
What is really happening when ``randomness'' produces stable averages? The mechanism
is straightforward. Each flip contributes a tiny, bounded nudge to the running fraction, and
the nudges are not coordinated -- they do not conspire to all push the same way.
With many independent nudges, the ups and downs largely cancel, and what is left is
the underlying tendency, exposed. The cancellation is the whole story. Later we
will make it precise as the Law of Large Numbers (the average converges) and the
Central Limit Theorem (the leftover wiggle has a universal bell shape). Keep this
image: \emph{randomness in the small, law in the large, with cancellation as the
bridge between them.}
\end{happeningbox}

\section{Randomness vs.\ Determinism}
\label{ch01:randomness}

Here is a question that troubles careful people, and it should: is the coin
\emph{really} random, or does it just look that way?

Press on the physics. A coin is a rigid body obeying Newton's laws. If you knew
its exact initial position, spin, velocity, the air currents, the elasticity of the
table -- in principle you could compute the outcome exactly. There is no genuine
chance in there at all. The same is true of dice, of shuffled cards, of the
roulette wheel. These are \emph{deterministic} systems. So in what sense are they
random?

\begin{focusbox}
The resolution is that probability is, very often, a confession of
\emph{ignorance}, not a claim about the universe's deepest nature. The coin's
outcome is determined; it is determined by quantities we do not and practically
cannot measure. The system is so sensitive to its starting conditions that a
difference too small to see -- a hair's width in launch angle -- flips the result.
This is the hallmark of \emph{chaos}: deterministic dynamics that amplify tiny
unknowns into total unpredictability. Probability steps in precisely where
determinism becomes useless to us. The coin is not random; \emph{we} are
ignorant, in a sharp and quantifiable way, and probability is the mathematics of
that ignorance.
\end{focusbox}

This is liberating, because it means we do not have to settle the metaphysics of
free will and quantum mechanics before we can do statistics. Whether or not the
universe is ``really'' random at bottom (quantum mechanics suggests some of it
genuinely is), the \emph{practical} situation -- a system whose relevant details we
cannot pin down -- is identical, and the same mathematics applies.

\begin{connectionbox}
There is a deep link to dynamical systems and chaos theory here, and through it to
calculus and differential equations. A chaotic map like the ``doubling map''
$x \mapsto 2x \bmod 1$ takes a perfectly deterministic rule and, by stretching and
folding the interval, produces sequences that pass every statistical test for
randomness. Sensitive dependence on initial conditions -- the ``butterfly effect''
-- is exactly the engine that converts unmeasurable precision into effective
chance. If you want to see how a single deterministic equation can manufacture
apparent randomness, that is the doorway, and it connects this book to the calculus
companion's treatment of dynamical systems.
\end{connectionbox}

\begin{paradoxbox}
\textbf{Is a fixed-but-unknown digit ``random''?} Consider the trillionth digit of
$\pi$. It is a specific, fixed number -- there is nothing chancy about it; $\pi$
does not change. Yet if you have not looked it up, you would happily assign each
digit $0$ through $9$ a probability of $1/10$, and you would be \emph{right} to, in
the only sense that matters for betting. This is the cleanest possible case where
frequentist language fails (the digit is not the outcome of any repeatable
experiment) and Bayesian language thrives (it is a perfect picture of degree of
belief under ignorance). The ``paradox'' -- assigning probabilities to a thing with
no randomness in it -- dissolves the moment you accept that probability can measure
\emph{your information}, not just the world's caprice. Many real problems, from
cryptography to forecasting, are exactly of this type.
\end{paradoxbox}

\section{A First Pass at the Rules}
\label{ch01:rules}

Enough philosophy; let us start computing. We will keep it light and concrete
here. We will develop the essential counting tools in Chapter~2, and build the
formal scaffolding (sample spaces, events, Kolmogorov's axioms) in Chapter~3.
The goal now is to get the \emph{feel} of how probabilities behave, using
examples a child could check.

\subsection{Outcomes, events, and the basic vocabulary}
\label{ch01:vocab}

Roll a single fair die. The possible \textbf{outcomes} are $1,2,3,4,5,6$. The
complete list of outcomes is the \textbf{sample space}, written $\Omega$ (capital
omega). An \textbf{event} is any collection of outcomes -- a question you can ask
that the experiment answers yes or no. ``The roll is even'' is the event
$\set{2,4,6}$. ``The roll is at least $5$'' is $\set{5,6}$.

\begin{definition}[Equally likely outcomes]
\label{ch01:def-equally-likely}
If a sample space $\Omega$ has finitely many outcomes and there is no reason to
favor any one over another -- by symmetry, as with a fair die or a well-shuffled
deck -- then we assign each outcome the same probability, and for any event $A$,
\[
  \P(A) \;=\; \frac{\text{number of outcomes in } A}{\text{number of outcomes in } \Omega}
  \;=\; \frac{\abs{A}}{\abs{\Omega}},
\]
where $\abs{A}$ denotes the number of outcomes in the event $A$, and
$\abs{\Omega}$ the total number in the sample space. Probability becomes
counting: favorable cases over total cases.
\end{definition}

This is the oldest definition in the book -- the ``classical'' definition of
Laplace -- and it carries us a long way. The probability the die shows an even
number is $\abs{\set{2,4,6}} / \abs{\set{1,\dots,6}} = 3/6 = 1/2$. The probability
it is at least $5$ is $2/6 = 1/3$.

\begin{trapbox}
The phrase ``no reason to favor any one'' is doing heavy lifting, and it can betray
you. It is tempting to say ``there are two outcomes, rain or no rain, so each has
probability $1/2$.'' That is nonsense -- the two are not symmetric. Equal
likelihood must be \emph{justified}, by a physical symmetry or by genuine
ignorance that is itself symmetric, not assumed for free just because you can list
two cases. We will meet a sharp version of this trap, Bertrand's paradox, later on.
\end{trapbox}

\subsection{Three rules you already believe}
\label{ch01:three-rules}

Even without axioms, you already accept three things about probabilities. Let us
name them, because everything else is built from them.

\textbf{Rule 1 (Bounds).} A probability is a number between $0$ and $1$. An
impossible event has probability $0$; a certain event has probability $1$. You
cannot be ``$140\%$ sure'' and you cannot be ``$-20\%$ sure.'' In symbols,
$0 \le \P(A) \le 1$, with $\P(\Omega) = 1$.

\textbf{Rule 2 (Addition for incompatible events).} If two events cannot both
happen -- they are \emph{mutually exclusive}, or \emph{disjoint} -- then the
probability that one or the other happens is the sum of their probabilities. The
die cannot be both a $2$ and a $5$, so
\[
  \P(\text{roll is } 2 \text{ or } 5) = \P(2) + \P(5) = \tfrac16 + \tfrac16 = \tfrac13.
\]

\textbf{Rule 3 (Complement).} The probability that $A$ does \emph{not} happen is
$1$ minus the probability that it does: $\P(\text{not } A) = 1 - \P(A)$. This
follows from Rules 1 and 2, since $A$ and ``not $A$'' are incompatible and together
cover everything. It is the single most useful trick in elementary probability,
because ``not'' is often far easier to count.

\begin{whybox}
Why does the complement rule follow? The events $A$ and ``not $A$'' are disjoint
(an outcome cannot be both in $A$ and not in $A$) and exhaustive (every outcome is
in one of them). So by Rule 2 their probabilities add, and by Rule 1 that total is
$\P(\Omega) = 1$:
\[
  \P(A) + \P(\text{not } A) = \P(\Omega) = 1
  \quad\Longrightarrow\quad
  \P(\text{not } A) = 1 - \P(A).
\]
Notice we proved a rule from simpler rules. That is the texture of the whole
subject: a tiny set of obvious-sounding axioms, and a tower of consequences. In
this section you saw the words; in Chapter~3 we make these the official axioms
(due to Andrey Kolmogorov in 1933) and watch the tower rise.
\end{whybox}

Figure~\ref{ch01:fig-probability-area} brings these three rules into a single
geometric picture: probability as area.

\begin{figure}[H]
\centering
\begin{tikzpicture}[scale=0.92]
  \begin{scope}[shift={(0,0)}]
    \draw[draw=gray!50, fill=gray!5, thick, rounded corners=3pt] (0,0) rectangle (6.2, 3.8);
    \node[font=\footnotesize\bfseries, color=gray!70!black, anchor=north east] at (6.0, 3.6) {Sample space $\Omega$ (Area $= 1$)};

    \draw[draw=sourceblue!80, fill=sourceblue!25, thick] (1.9, 1.8) circle (1.15);
    \node[font=\small\bfseries, color=sourceblue!90!black] at (1.9, 2.05) {$A$};
    \node[font=\footnotesize, color=sourceblue!90!black] at (1.9, 1.55) {$\P(A)$};

    \draw[draw=quotientorange!80, fill=quotientorange!25, thick] (4.6, 1.8) circle (0.9);
    \node[font=\small\bfseries, color=quotientorange!90!black] at (4.6, 2.05) {$B$};
    \node[font=\footnotesize, color=quotientorange!90!black] at (4.6, 1.55) {$\P(B)$};

    \node[font=\footnotesize, align=center] at (3.1, -0.45) {
      \textbf{Disjoint events} ($A \cap B = \emptyset$):\\
      $\P(A \text{ or } B) = \P(A) + \P(B)$
    };
  \end{scope}

  \begin{scope}[shift={(7.3,0)}]
    \draw[draw=gray!50, fill=gray!15, thick, rounded corners=3pt] (0,0) rectangle (6.2, 3.8);
    \node[font=\footnotesize\bfseries, color=gray!70!black, anchor=north east] at (6.0, 3.6) {Sample space $\Omega$ (Area $= 1$)};

    \draw[draw=sourceblue!80, fill=sourceblue!30, thick] (2.2, 1.8) circle (1.15);
    \node[font=\small\bfseries, color=sourceblue!90!black] at (2.2, 2.05) {$A$};
    \node[font=\footnotesize, color=sourceblue!90!black] at (2.2, 1.55) {$\P(A)$};

    \node[font=\small\bfseries, color=gray!80!black] at (4.7, 2.05) {$\text{not } A$};
    \node[font=\footnotesize, color=gray!80!black] at (4.7, 1.55) {$1 - \P(A)$};

    \node[font=\footnotesize, align=center] at (3.1, -0.45) {
      \textbf{Complement rule}:\\
      $\P(\text{not } A) = 1 - \P(A)$
    };
  \end{scope}
\end{tikzpicture}
\caption{Probability as area in the sample space $\Omega$. The entire canvas has total area $\P(\Omega) = 1$ (Rule~1). Left: when events $A$ and $B$ cannot co-occur (they are disjoint), their patches do not overlap, so the probability of either happening is simply the sum of their individual areas (Rule~2). Right: an event $A$ and its complement ``$\text{not } A$'' together tile the entire space without overlap, forcing $\P(\text{not } A) = 1 - \P(A)$ (Rule~3).}
\label{ch01:fig-probability-area}
\end{figure}

\begin{example}[The power of looking at the complement]
\label{ch01:ex-complement}
Roll a fair die four times. What is the probability of getting \emph{at least one}
six?

Counting ``at least one'' directly is a headache: exactly one six, exactly two,
exactly three, exactly four -- four cases to tally and add. But the complement of
``at least one six'' is ``\emph{no} sixes at all,'' which is one clean case. Each
roll avoids six with probability $5/6$, and (because the rolls do not influence
each other) the chance all four avoid six is $(5/6)^4$. Therefore
\[
  \P(\text{at least one six}) = 1 - \left(\tfrac{5}{6}\right)^4
  = 1 - \tfrac{625}{1296} = \tfrac{671}{1296} \approx 0.518.
\]
Just over an even chance. We will formalize the multiplication rule and
independence in Chapters~3 and~4; here, savor that flipping to the complement
turned four cases into one.
\end{example}

\begin{walkbox}
\textbf{Pause and think: why the ``disjoint'' condition matters.}
Roll a single fair die. Let $A$ be the event ``the roll is odd'' ($\P(A) = 1/2$)
and $B$ be the event ``the roll is at least $4$'' ($\P(B) = 1/2$). What is the
probability that the roll is odd \emph{or} at least $4$? Can you simply add
$1/2 + 1/2 = 1$?

Check the outcomes: the roll $5$ is both odd \emph{and} at least $4$. Because
$A$ and $B$ overlap, adding their probabilities counts outcome $5$ twice. The
actual outcomes in ``$A$ or $B$'' are $\set{1, 3, 4, 5, 6}$, five of the six
equally likely faces, so the true probability is $5/6$, not $1$. Rule~2 requires
incompatible (disjoint) events for a reason. When events overlap, we must
subtract the double-counted intersection -- a rule we will make official as the
inclusion--exclusion formula in Chapter~3.
\end{walkbox}

\subsection{Why these rules and not others?}
\label{ch01:why-these}

It would be fair to ask whether we could have written down \emph{different} rules.
Why addition for disjoint events? Why a ceiling of $1$? The frequentist and
Bayesian each answer in their own voice, and -- as promised -- they land in the
same place.

\begin{whybox}
\emph{The frequentist's reason.} Probabilities are limiting frequencies, and
frequencies directly obey these rules at every finite stage, so they must in the
limit too. In $n$ flips, the fraction of heads is between $0$ and $1$ -- there is
Rule 1. If $A$ and $B$ never co-occur, then the count of ``$A$ or $B$'' is exactly
the count of $A$ plus the count of $B$, so the fractions add -- there is Rule 2.
The fraction of ``not heads'' is $1$ minus the fraction of heads -- there is
Rule 3. The axioms are nothing more than the bookkeeping of counting, carried to
the limit.

\emph{The Bayesian's reason.} Probabilities are coherent betting prices. If you
price bets in a way that violates these rules -- say you are willing to pay $0.6$
for a bet on $A$ \emph{and} $0.6$ for a bet on ``not $A$'' -- then a sharp operator
can sell you both bets, collect $1.20$, and pay out exactly $1.00$ no matter what
happens, pocketing $0.20$ from you with certainty. The only prices immune to such a
guaranteed loss (a ``Dutch book'') are exactly those satisfying Rules 1--3. Coherence
\emph{forces} the axioms.
\end{whybox}

That two such different stories -- counting in the limit, and refusing to be
swindled -- arrive at the identical algebra is the first sign that we have found
something real, not arbitrary. We are not inventing probability; we are
discovering the only consistent way to reason under uncertainty.

\subsection{A tiny worked tree}
\label{ch01:tree}

One more taste before the formalities. Probabilities chain together along
sequences of events, and a \emph{tree diagram} is the natural picture.
Figure~\ref{ch01:fig-tree} shows two flips of a fair coin.

\begin{figure}[H]
\centering
\begin{tikzpicture}[
  level distance=26mm,
  level 1/.style={sibling distance=34mm},
  level 2/.style={sibling distance=18mm},
  edge from parent/.style={draw, -{Latex[length=2mm]}, sourceblue, thick},
  every node/.style={font=\small},
]
\node[circle, draw=gray!60, fill=gray!8, minimum size=7mm] {}
  child {node[circle, draw=gray!60, fill=blue!8, minimum size=7mm] {H}
    child {node[draw, rounded corners, fill=green!8] {HH}
      edge from parent node[left] {$\frac12$}}
    child {node[draw, rounded corners, fill=green!8] {HT}
      edge from parent node[right] {$\frac12$}}
    edge from parent node[above left] {$\frac12$}}
  child {node[circle, draw=gray!60, fill=blue!8, minimum size=7mm] {T}
    child {node[draw, rounded corners, fill=green!8] {TH}
      edge from parent node[left] {$\frac12$}}
    child {node[draw, rounded corners, fill=green!8] {TT}
      edge from parent node[right] {$\frac12$}}
    edge from parent node[above right] {$\frac12$}};
\end{tikzpicture}
\caption{Two flips of a fair coin. Each path from root to leaf is one outcome; its
probability is the product of the edge probabilities along the way, here
$\frac12 \cdot \frac12 = \frac14$. The four leaves are equally likely and their
probabilities sum to $1$, as they must.}
\label{ch01:fig-tree}
\end{figure}

Reading the tree: multiply \emph{along} a branch (the two flips happen in
sequence), and add \emph{across} branches (different ways to succeed are disjoint).
What is the probability of exactly one head? Two leaves qualify, HT and TH, each
with probability $1/4$, and they are disjoint, so the answer is
$1/4 + 1/4 = 1/2$. Multiply along, add across -- you will use that rhythm for the
rest of your life.

\section{Where It All Came From: Games of Chance}
\label{ch01:history}

Mathematics usually grows from one of two roots: a practical problem that will not
go away, or a puzzle too pretty to ignore. Probability grew from both at once, and
the soil was the gambling table.

Its first true pioneer was Gerolamo Cardano (1501--1576): a Milanese physician,
astrologer, and inveterate gambler who solved cubic equations and gambled almost
daily for decades. Around 1564 Cardano wrote \emph{Liber de Ludo Aleae} (``The
Book on Games of Chance''), the first serious attempt to put gambling odds on a
mathematical footing. In it he stated the core intuition of the classical
definition: for fair dice, the probability of an event is the number of favorable
cases divided by the total number of equally likely cases. Cardano's key
breakthrough was realizing that when you roll two dice, there are $36$ equally
likely outcomes, not $11$ (the common mistake of counting only the
distinguishable sums), and he used this to price wagers correctly. His book sat
unpublished until 1663, long after his death, so its immediate influence was
muted -- but Cardano had glimpsed, in isolation, the foundational technique of
the entire subject: \emph{enumerate the equally likely elementary cases and
count}.

Cardano's insight about two dice is worth a picture, because it exposes the single
most common counting trap in elementary probability. The eleven possible
\emph{sums} ($2$ through $12$) are emphatically not equally likely: there is
only one way to roll a $2$ (namely $1{+}1$) but six ways to roll a $7$ ($1{+}6$,
$2{+}5$, $3{+}4$, $4{+}3$, $5{+}2$, $6{+}1$). The equally likely objects are the
$36$ \emph{ordered pairs}, and the probability of each sum is just the tally of
pairs along each anti-diagonal. Figure~\ref{ch01:fig-dice} shows the resulting
tent shape.

\begin{figure}[H]
\centering
\begin{tikzpicture}
  % Panel (a): 6x6 grid of outcomes
  \begin{scope}[shift={(0,0)}]
    \node[font=\footnotesize\bfseries, anchor=center] at (1.25, 4.3) {All $36$ equally likely pairs $(d_1, d_2)$};
    \node[font=\scriptsize\color{gray!70!black}, anchor=east] at (-0.65, 1.75) {Die 1 ($d_1$)};
    \node[font=\scriptsize\color{gray!70!black}, anchor=south] at (1.25, 3.65) {Die 2 ($d_2$)};
    
    \foreach \r in {1,...,6} {
      \node[font=\tiny\bfseries, color=gray!80!black] at (-0.35, {3.5 - 0.5*\r}) {\r};
      \node[font=\tiny\bfseries, color=gray!80!black] at ({0.5*\r - 0.25}, 3.45) {\r};
    }
    
    \foreach \r in {1,...,6} {
      \foreach \c in {1,...,6} {
        \pgfmathtruncatemacro{\s}{\r + \c}
        \ifnum\s=7
          \fill[sourceblue!30, rounded corners=1pt] ({0.5*\c - 0.48}, {3.5 - 0.5*\r - 0.23}) rectangle ++(0.46, 0.46);
          \draw[sourceblue!80, line width=0.6pt, rounded corners=1pt] ({0.5*\c - 0.48}, {3.5 - 0.5*\r - 0.23}) rectangle ++(0.46, 0.46);
          \node[font=\tiny\bfseries, color=sourceblue!90!black] at ({0.5*\c - 0.25}, {3.5 - 0.5*\r}) {\s};
        \else
          \ifnum\s=6
            \fill[targetgreen!20, rounded corners=1pt] ({0.5*\c - 0.48}, {3.5 - 0.5*\r - 0.23}) rectangle ++(0.46, 0.46);
            \draw[targetgreen!60, line width=0.4pt, rounded corners=1pt] ({0.5*\c - 0.48}, {3.5 - 0.5*\r - 0.23}) rectangle ++(0.46, 0.46);
            \node[font=\tiny, color=targetgreen!90!black] at ({0.5*\c - 0.25}, {3.5 - 0.5*\r}) {\s};
          \else
            \ifnum\s=8
              \fill[targetgreen!20, rounded corners=1pt] ({0.5*\c - 0.48}, {3.5 - 0.5*\r - 0.23}) rectangle ++(0.46, 0.46);
              \draw[targetgreen!60, line width=0.4pt, rounded corners=1pt] ({0.5*\c - 0.48}, {3.5 - 0.5*\r - 0.23}) rectangle ++(0.46, 0.46);
              \node[font=\tiny, color=targetgreen!90!black] at ({0.5*\c - 0.25}, {3.5 - 0.5*\r}) {\s};
            \else
              \fill[gray!6, rounded corners=1pt] ({0.5*\c - 0.48}, {3.5 - 0.5*\r - 0.23}) rectangle ++(0.46, 0.46);
              \draw[gray!35, line width=0.3pt, rounded corners=1pt] ({0.5*\c - 0.48}, {3.5 - 0.5*\r - 0.23}) rectangle ++(0.46, 0.46);
              \node[font=\tiny, color=gray!60!black] at ({0.5*\c - 0.25}, {3.5 - 0.5*\r}) {\s};
            \fi
          \fi
        \fi
      }
    }
    \node[font=\scriptsize, align=left, color=sourceblue!80!black] at (1.5, -0.3) {Longest diagonal: sum $= 7$ ($6$ pairs)};
  \end{scope}

  % Arrow
  \draw[->, >=stealth, thick, gray!60] (3.1, 1.7) -- (3.9, 1.7);
  \node[above=2pt, font=\tiny\bfseries, text=gray!70!black] at (3.5, 1.7) {count};
  \node[below=2pt, font=\tiny, text=gray!60!black] at (3.5, 1.7) {diagonals};

  % Panel (b): Bar plot
  \begin{scope}[shift={(6.1,-0.5)}]
    \begin{axis}[
      width=7.0cm, height=5.2cm,
      title style={font=\small},
      title={Probability distribution of the sum $d_1 + d_2$},
      ybar, bar width=8pt,
      xlabel={sum $s$}, ylabel={probability},
      xmin=1.4, xmax=12.6, ymin=0, ymax=0.22,
      xtick={2,3,4,5,6,7,8,9,10,11,12},
      ytick={0,0.05,0.10,0.15},
      grid=both, grid style={gray!18},
      nodes near coords, every node near coord/.append style={font=\tiny, color=gray!55!black, yshift=4pt},
      point meta=explicit symbolic,
    ]
    \addplot[draw=distblue, fill=histfill] coordinates {
      (2,0.0278)[$\frac{1}{36}$] (3,0.0556)[$\frac{2}{36}$] (4,0.0833)[$\frac{3}{36}$]
      (5,0.1111)[$\frac{4}{36}$] (6,0.1389)[$\frac{5}{36}$] (7,0.1667)[$\frac{6}{36}$]
      (8,0.1389)[$\frac{5}{36}$] (9,0.1111)[$\frac{4}{36}$] (10,0.0833)[$\frac{3}{36}$]
      (11,0.0556)[$\frac{2}{36}$] (12,0.0278)[$\frac{1}{36}$]
    };
    \end{axis}
  \end{scope}
\end{tikzpicture}
\caption{The geometry behind two dice. Left: the $36$ equally likely elementary outcomes $(d_1, d_2)$ form a square grid. The sum $d_1+d_2$ is constant along anti-diagonals, with the longest diagonal containing the $6$ pairs that sum to $7$ (blue). Right: collapsing each diagonal into a single bar produces the symmetric tent distribution. Treating the eleven sums as equally likely would flatten this tent into a horizontal line.}
\label{ch01:fig-dice}
\end{figure}

While Cardano worked in private, the subject began to develop publicly in the summer
of 1654, in one of the most consequential exchanges of letters in the history of
science.

\begin{historybox}
\textbf{The Chevalier's dice, the problem of points, and the birth of probability (1654).}
In the summer of 1654, an amateur mathematician and high-stakes Parisian gambler,
Antoine Gombaud (who styled himself the Chevalier de M\'er\'e), brought two
baffling puzzles to Blaise Pascal.

The first was a practical dice puzzle that seemed to defy reason. Gamblers knew
that betting on at least one six in $4$ rolls of a single die gave a slight
advantage (as in Example~\ref{ch01:ex-complement}, $\P \approx 0.518 > 1/2$).
De M\'er\'e reasoned by naive proportion: since two dice produce $36$ outcomes
instead of $6$, rolling two dice $24$ times should give the exact same edge
($24/36 = 4/6 = 2/3$). Yet in practice, betting on at least one double-six in
$24$ rolls steadily drained his purse. De M\'er\'e was so perturbed that he
declared arithmetic itself was defective.

The second puzzle was an old and notorious legal and sporting question: the
\emph{problem of points}. Two players in a fair game of chance (say, first to
$3$ wins) are interrupted mid-play with the score at $2$--$1$. How should the
prize stakes be split fairly? Splitting $2{:}1$ by the current score rewards past
victories while ignoring how few rounds remain; splitting equally ignores that
one player is one point from victory.

Working with only pen, paper, and basic arithmetic---with no symbolic notation
for probability (symbols like $\P(A)$ were still centuries away) and no
calculus---Pascal initiated a correspondence with Pierre de Fermat in Toulouse.
Over a handful of letters, they laid the foundations of mathematical probability.
Yet they solved the problem of points by two entirely different methods, sparking
a famous debate. Fermat used combinatorial enumeration: he imagined the game
extended by ``phantom rounds'' until its maximum possible duration was reached,
ensuring all branches in the tree of possibilities had equal length and equal
weight. Pascal initially objected that playing phantom rounds after someone had
already won was unnatural fiction, preferring a step-by-step backward recursion
from final payouts (which introduced the concept of mathematical expectation).
Fermat proved that phantom rounds were mathematically exact: they preserve the
correct weights while making every branch equally likely.

Pascal also resolved de M\'er\'e's dice puzzle: the probability of avoiding a
double-six in one roll of two dice is $35/36$, so over $24$ rolls the chance of
\emph{never} seeing a double-six is $(35/36)^{24} \approx 0.5086$. The probability
of winning is therefore $1 - 0.5086 = 0.4914$---just below an even chance ($49.1\%$).
Arithmetic was not broken; de M\'er\'e's naive linear proportion was simply
invalid because probabilities of non-disjoint events do not add linearly.

This correspondence transformed chance from mystical fortune into a deductive
science. Dutch physicist Christiaan Huygens learned of the letters and in 1657
published the first probability textbook (\emph{De ratiociniis in ludo aleae}),
and within decades Jacob Bernoulli proved the first Law of Large Numbers. What
Pascal and Fermat left unresolved, however, was the ultimate meaning of the
numbers they computed: does probability describe the physical frequency of an
unfolding process, or a thinking agent's rational degree of belief under
imperfect information? That question remains an active and vigorous debate in
modern science.
\end{historybox}

Let us actually solve a clean version of their problem, because it is short, it is
beautiful, and it shows the new style of thinking in action.

\begin{example}[A small problem of points]
\label{ch01:ex-points}
Alice and Bob play rounds of a fair game; each round, either wins with probability
$1/2$, independently. The first to win $3$ rounds takes the whole pot of \$64.
The game is interrupted when Alice has $2$ wins and Bob has $1$. How should the
\$64 be split fairly?

The wrong-but-tempting answer is $2{:}1$, splitting by current score. The right
answer asks: if they \emph{had} continued, how likely was each to win? Alice needs
just $1$ more round; Bob needs $2$.

At most $2$ more rounds settle it. Enumerate the four equally likely two-round
continuations -- writing A for ``Alice wins the round'' and B for ``Bob wins the
round'':
\[
  \text{AA},\quad \text{AB},\quad \text{BA},\quad \text{BB},
  \qquad \text{each with probability } \tfrac14.
\]
Alice wins the pot in any continuation containing at least one A among the rounds
she needs -- that is, in AA, AB, and BA (in each, she reaches $3$ wins). Bob wins
only in BB, where he takes both remaining rounds to reach $3$. So
\[
  \P(\text{Alice wins}) = \tfrac34, \qquad \P(\text{Bob wins}) = \tfrac14.
\]
Fair division: Alice gets $\frac34 \cdot \$64 = \$48$, Bob gets
$\frac14 \cdot \$64 = \$16$. Notice the split is $3{:}1$, not the naive $2{:}1$ --
being one round from victory is worth far more than the score alone suggests. This
is the entire Pascal--Fermat method: enumerate the equally likely futures, weight
them, and combine.
\end{example}

\begin{figure}[H]
\centering
\begin{tikzpicture}[
  state/.style={rectangle, draw=sourceblue!70, fill=sourceblue!8, rounded corners=2pt, font=\small, minimum height=0.7cm, minimum width=2.4cm, inner sep=4pt, align=center},
  winA/.style={rectangle, draw=targetgreen!80, fill=targetgreen!12, rounded corners=2pt, font=\small\bfseries, minimum height=0.7cm, minimum width=2.4cm, inner sep=4pt, align=center},
  winB/.style={rectangle, draw=quotientorange!80, fill=quotientorange!12, rounded corners=2pt, font=\small\bfseries, minimum height=0.7cm, minimum width=2.4cm, inner sep=4pt, align=center},
  trans/.style={->, >=stealth, thick, draw=gray!70!black},
  lbl/.style={font=\scriptsize, fill=white, inner sep=1.5pt}
]
  \node[state] (S0) at (0, 1.2) {Current State\\Alice $2$, Bob $1$};
  \node[winA] (A1) at (4.5, 2.4) {Alice wins\\(pot $\$64$)};
  \node[state] (S1) at (4.5, 0.0) {Alice $2$, Bob $2$};
  \node[winA] (A2) at (9.0, 0.9) {Alice wins\\(pot $\$64$)};
  \node[winB] (B1) at (9.0, -0.9) {Bob wins\\(pot $\$64$)};

  \draw[trans] (S0) -- node[lbl, sloped, above=1pt] {Alice ($1/2$)} (A1);
  \draw[trans] (S0) -- node[lbl, sloped, below=1pt] {Bob ($1/2$)} (S1);
  \draw[trans] (S1) -- node[lbl, sloped, above=1pt] {Alice ($1/2$)} (A2);
  \draw[trans] (S1) -- node[lbl, sloped, below=1pt] {Bob ($1/2$)} (B1);

  \node[font=\footnotesize, right=0.2cm of A1, color=targetgreen!80!black] {Prob: $\frac12$};
  \node[font=\footnotesize, right=0.2cm of A2, color=targetgreen!80!black] {Prob: $\frac12 \cdot \frac12 = \frac14$};
  \node[font=\footnotesize, right=0.2cm of B1, color=quotientorange!80!black] {Prob: $\frac12 \cdot \frac12 = \frac14$};

  \node[font=\small, align=center] at (4.5, -1.8) {
    Alice: $\P(\text{wins}) = \frac12 + \frac14 = \frac34 \implies \text{fair share } \$48$ \\[5pt]
    Bob: $\P(\text{wins}) = \frac14 \implies \text{fair share } \$16$
  };
\end{tikzpicture}
\caption{The decision tree for the Problem of Points (Example~\ref{ch01:ex-points}). From the interrupted state, Alice needs one win while Bob needs two. Branching round by round shows Alice wins either on round $1$ (prob $\frac12$) or round $2$ (prob $\frac14$), yielding total probability $\frac34$.}
\label{ch01:fig-points-tree}
\end{figure}

\begin{remark}
A subtlety worth noticing: we pretended both remaining rounds are played even
though the game would actually stop the instant someone reaches $3$. This is
harmless -- playing ``phantom'' extra rounds does not change who has already won,
and it makes every leaf equally likely so we can just count. Pascal and Fermat
argued about exactly this point in their letters; Fermat favored the enumeration
of all continuations, Pascal a recursive ``expected value'' bookkeeping. Both are
correct and they agree. We will meet both viewpoints again, formalized, when we
study expectation.
\end{remark}

\begin{applicationbox}
\textbf{From the card table to the modern world.} It is easy to smile at the idea
of probability being born to settle gambling debts, but the through-line to today
is direct and unbroken. The ``problem of points'' is, structurally, the problem an
insurance company solves when it prices a policy on an uncertain future, the
problem a finance desk solves when it values an option whose payoff depends on
where a stock lands, and the problem a clinical trial faces when it decides whether
to stop early because one treatment is pulling ahead. Each asks: \emph{given where
we are now, weight the possible futures and combine them into a fair present
value.} Every time your phone's keyboard predicts the next word, every time a spam
filter scores an email, every time a navigation app estimates your arrival time,
the machine is enumerating possibilities, weighting them by probability, and
combining -- the Pascal--Fermat move, executed billions of times a second. The
card table was just the first place humans were forced to get it right.
\end{applicationbox}

\section{The Understanding-First Promise}
\label{ch01:promise}

Before the exercises, a word about how to read the rest of this book, since this
chapter sets the tone for everything after.

\begin{focusbox}
Probability is unusually punishing to memorize and unusually rewarding to
understand. Memorize the formula for conditional probability and you will misapply
it within a week; understand the tree picture behind it and you will reconstruct
the formula whenever you need it, forever. So throughout this book we will always
do the same three things, in this order. \emph{First}, build the idea from a
concrete story until it feels obvious. \emph{Second}, make it precise -- a real
definition, a real derivation, no hand-waving disguised as rigor. \emph{Third},
stress-test it against a paradox, a real application, or a connection to other
mathematics, putting the theory to work. If at any point a formula feels like a
routine you are executing mechanically, pause -- you may have skipped step one,
and the best remedy is to return to a small, concrete example you can work by hand.
\end{focusbox}

\begin{recoverybox}
\textbf{Where we have been.} Probability is a number, between $0$ and $1$, that
measures uncertainty. It has two solid interpretations -- long-run frequency
(natural for repeatable experiments) and degree of belief (natural for one-off
questions) -- and, remarkably, both obey the \emph{same} three rules: probabilities
lie in $[0,1]$, disjoint events add, and complements subtract from $1$. We need
mathematics because intuition fails systematically at chance, because uncertainties
must be combined consistently, and because stable laws hide inside aggregate
randomness. ``Randomness'' is often a confession of ignorance about a deterministic
but chaotic system. And the whole subject was born at the gambling table, in
Cardano's systematic counting and the Pascal--Fermat letters, from a single move:
enumerate the equally likely cases, weight them, combine.

\textbf{Where we are going next.} To use the classical formula $\P(A) = \abs{A}/\abs{\Omega}$, you have to count. For a single die or two coins, counting outcomes is easy. But what if you are dealt a five-card poker hand from a deck of $52$, or want to know the chance that two people in a room share a birthday? Listing outcomes by hand is impossible. Before we make $\Omega$, events, and Kolmogorov's axioms official in Chapter~3, we need the mathematical art of counting large structured sets without listing them one by one: combinatorics. That is exactly what Chapter~2 builds from the ground up.
\end{recoverybox}

\section{Exercises}
\label{ch01:exercises}

\begin{exercise}[Warm-up mental checks]
\label{ch01:ex-mental}
Answer each mentally before writing anything down:
\begin{enumerate}
\item You flip a fair coin $3$ times. What is the probability of seeing at least
one head? Compute it in under five seconds using Rule~3.
\item If an event $A$ satisfies $\P(A) = 0.75$, what is $\P(\text{not } A)$?
\item Can two mutually exclusive events $A$ and $B$ have probabilities $\P(A) = 0.6$
and $\P(B) = 0.5$? Explain why or why not using Rules~1 and~2.
\item Two fair dice are rolled. Which sum is the single most likely outcome, and
what is its exact probability?
\end{enumerate}
\end{exercise}

\begin{exercise}[Frequency in your hand]
\label{ch01:ex-sim}
Flip a real coin $20$ times and record the running fraction of heads after each
flip. Plot it (by hand is fine) against the flip number. Does it wander early and
settle late, as in Figure~\ref{ch01:fig-stabilize}? Now imagine doing this for
$2000$ flips: describe in one sentence what the plot would look like, and say
precisely which interpretation of probability -- frequentist or Bayesian -- this
exercise is illustrating.
\end{exercise}

\begin{exercise}[Naming the interpretation]
\label{ch01:ex-name}
For each statement, say whether it is most naturally frequentist, most naturally
Bayesian, or genuinely either, and explain in one sentence why.
\begin{enumerate}
\item ``The probability this fair die shows a $4$ is $1/6$.''
\item ``There is a $30\%$ chance it rains in Lagos tomorrow.''
\item ``The probability the defendant is guilty, given the evidence, is high.''
\item ``About $0.7\%$ of manufactured chips from this line are defective.''
\item ``The probability that the millionth digit of $\pi$ is a $7$ is $1/10$.''
\end{enumerate}
\end{exercise}

\begin{exercise}[Complement practice]
\label{ch01:ex-comp}
Use the complement rule to answer each, and in each case say why ``not'' is easier
than direct counting.
\begin{enumerate}
\item Roll a fair die three times. Probability of at least one even number?
\item Flip a fair coin five times. Probability of at least one head?
\item Roll two fair dice. Probability the sum is \emph{not} $7$? (There are $36$
equally likely outcomes; first count the ways to get $7$.)
\end{enumerate}
\end{exercise}

\begin{exercise}[The Chevalier de M\'er\'e's dice puzzle]
\label{ch01:ex-demere}
Recall the historical betting puzzle that launched Pascal and Fermat's
correspondence:
\begin{enumerate}
\item In a single roll of a fair die, what is the probability of avoiding a six?
Roll four times: use the complement rule to show that the probability of getting
at least one six is $1 - (5/6)^4 \approx 0.5177$.
\item In a single roll of two fair dice, what is the probability of avoiding a
double-six? Roll twenty-four times: show that the probability of \emph{never}
seeing a double-six is $(35/36)^{24} \approx 0.5086$.
\item Deduce the probability of winning de M\'er\'e's wager (at least one
double-six in $24$ rolls). Explain in one or two sentences why de M\'er\'e's
linear scaling ($4/6 = 24/36$) led him to lose money in the long run.
\end{enumerate}
\end{exercise}

\begin{exercise}[A bigger problem of points]
\label{ch01:ex-points2}
Alice and Bob play the fair game of Example~\ref{ch01:ex-points} for a pot of
\$128, first to $4$ wins. The game is interrupted with Alice on $3$ wins and Bob on
$1$. At most how many more rounds settle it? Enumerate the equally likely
continuations, find each player's probability of winning, and give the fair split.
(Hint: as in the example, you may play phantom rounds so every leaf is equally
likely.)
\end{exercise}

\begin{exercise}[Spotting the false symmetry]
\label{ch01:ex-false-sym}
A friend argues: ``A tossed thumbtack lands either point-up or point-down. Two
outcomes, so the probability of point-up is $1/2$.'' Explain carefully what is
wrong with this reasoning, referencing Definition~\ref{ch01:def-equally-likely}.
How could you \emph{find} the true probability, and which interpretation of
probability does your method rely on?
\end{exercise}

\begin{exercise}[The Dutch book, made concrete]
\label{ch01:ex-dutch}
A bookmaker offers to sell you, for $\$0.60$ each, two tickets: ticket $A$ pays
$\$1$ if a certain horse wins, ticket $A'$ pays $\$1$ if it does not. You buy both.
\begin{enumerate}
\item How much do you pay in total, and how much do you collect, no matter what
happens?
\item Show that you lose a guaranteed amount, and state it.
\item Which of Rules 1--3 has the bookmaker's pricing violated? What is the largest
price he could charge for \emph{each} ticket without exposing himself to a sure
loss in the other direction?
\end{enumerate}
\end{exercise}

\begin{exercise}[Randomness from determinism]
\label{ch01:ex-doubling}
Consider the deterministic ``doubling map'': start with a number $x_0$ in $[0,1)$
and repeatedly apply $x_{n+1} = 2x_n \bmod 1$ (double it, drop the integer part).
Write each $x_n$ in binary. Show that applying the map simply \emph{shifts the
binary digits left by one}. Use this to explain, in a few sentences, why the
sequence of ``is $x_n \ge 1/2$?'' answers looks like fair coin flips even though the
rule is completely deterministic. What role does our ignorance of $x_0$'s exact
digits play? Connect this to the discussion in Section~\ref{ch01:randomness}.
\end{exercise}

\begin{exercise}[Two dice, the right count]
\label{ch01:ex-two-dice}
Cardano's contemporaries sometimes treated the sum of two dice as having $11$
equally likely values ($2$ through $12$). Explain why this is wrong by listing the
$36$ equally likely outcomes of an ordered pair of dice. Compute $\P(\text{sum}=7)$
and $\P(\text{sum}=2)$, and explain in one sentence the conceptual error in the
``$11$ equally likely sums'' view. (This is precisely the kind of mistake Cardano
avoided and that the classical definition is designed to prevent.)
\end{exercise}

\begin{exercise}[Conceptual: why infinity is uncomfortable]
\label{ch01:ex-infinity}
The frequentist defines $\P(\text{heads})$ as the limit of the fraction of heads as
the number of flips goes to infinity. Identify two distinct difficulties with
taking this as a \emph{definition} (as opposed to an interpretation). One should
concern what we can actually observe; the other should concern a hidden circularity
in the word ``converges.'' (You are not expected to resolve them here -- only to
articulate them clearly. The book resolves them when it reaches the Law of Large
Numbers.)
\end{exercise}

\begin{exercise}[Proof: the complement rule, twice]
\label{ch01:ex-proof}
Complete the following three steps:
\begin{enumerate}
\item Using only Rules 1 and 2 from Section~\ref{ch01:three-rules}, prove that
$\P(\text{not } A) = 1 - \P(A)$. State explicitly where you use that $A$ and
``not $A$'' are disjoint and where you use that together they are everything.
\item Deduce that the impossible event has probability $0$. (Hint: take $A = \Omega$.)
\item Deduce that if $A$ is contained in $B$ -- every outcome in $A$ is also in $B$
-- then $\P(A) \le \P(B)$. (Hint: write $B$ as the disjoint union of $A$ and ``the
part of $B$ outside $A$,'' then use Rule 1 on the leftover piece.)
\end{enumerate}
\end{exercise}
